% Este apartado está dado por el Sr. Bladimir Barron, ante cualquier duda es mejor dirigirse a él directamente.

% ======== Default Style ======== %

\lstset{
  inputencoding=utf8,
  extendedchars=true,
  basicstyle=\ttfamily\small,
  numbers=left,
  numberstyle=\tiny,
  stepnumber=1,
  numbersep=5pt,
  showstringspaces=false,
  frame=single,
  breaklines=true,
  tabsize=2,
  columns=fixed,
  keepspaces=true
}

% ======== Character Mapping ======== %
\lstset{ literate= {á}{{\'a}}1 {é}{{\'e}}1 {í}{{\'i}}1 {ó}{{\'o}}1 {ú}{{\'u}}1 {Á}{{\'A}}1 {É}{{\'E}}1 {Í}{{\'I}}1 {Ó}{{\'O}}1 {Ú}{{\'U}}1 {ñ}{{\~n}}1 {Ñ}{{\~N}}1 {ü}{{\"u}}1 {Ü}{{\"U}}1 {¿}{{\textquestiondown}}1 {¡}{{\textexclamdown}}1 }

% ======== TColorBox for listings ======== %

% Basic style for listings inside tcolorbox.
% It can be overwritten in the definitions being used 
\tcbset{ basicCodeStyle/.style={
  breakable, enhanced, listing only, fonttitle=\bfseries, sharp corners,
  boxrule=0.5pt, left=0mm, top=-1mm, bottom=2mm, overlay
  first={
    \node[anchor=south east, font=\itshape\small, text=black!85] at (frame.south
    east) {(Continúa $\rightarrow$) \,};
  },
  overlay
  middle={
    \node[anchor=south east, font=\itshape\small, text=black!85] at (frame.south
    east) {(Continúa $\rightarrow$) \,};
  },
  codetitle/.code={
    \tcbset{
      title={Código \thetcbcounter: ##1},
      title after break={Código \thetcbcounter: ##1 (continuación)}
    }
  }
} }

% ================================= %
% ======== Language Styles ======== %
% ================================= %

% ======== Python Style ======== %
\lstdefinestyle{pythonstyle}{
  language=Python, inputencoding=utf8, basicstyle=\ttfamily\small,
  keywordstyle=\color{blue}\bfseries, stringstyle=\color{teal},
  commentstyle=\color{black!70}\itshape, numbers=left, numberstyle=\tiny,
  frame=none, breaklines=true, tabsize=2, columns=fixed, keepspaces=true,
}

% ======== Insert Python ======== %
\newtcbinputlisting[auto counter, number within=noseccion]{\pycode}[2][]{
  basicCodeStyle, listing file={#2}, listing options={style=pythonstyle}, #1
}

% ======== Example of use 
% \pycode[codetitle={python code}, label=lst:python_example]{./path/to/python}

% ======== Fortran90 Style ======== %
\lstdefinestyle{fortranstyle}{
  language=[90]Fortran, basicstyle=\ttfamily\small,
  keywordstyle=\color{blue}\bfseries, commentstyle=\color{gray}\itshape,
  stringstyle=\color{teal}, numbers=left, numberstyle=\tiny, frame=none,
  breaklines=true, keepspaces=true, tabsize=2
}

% ======== Insert Fortran ======== %
\newtcbinputlisting[auto counter, number within=noseccion]{\fcode}[2][]{
  basicCodeStyle, listing file={#2}, listing options={style=fortranstyle}, #1
}

% ======== Example of use 
% \fcode[codetitle={F90 code}, label=lst:fortran_example]{./path/to/fortran}

% ======== Terminal Style ======== %
\lstdefinestyle{terminalstyle}{
  frame=none, numbers=none, showstringspaces=false, breaklines=true,
  keepspaces=true, columns=fixed,
}

% ======== Insert Code Output ======== %
\newtcbinputlisting[auto counter]{\codeoutput}[3][]{
  basicCodeStyle,
  bottom=0mm,
  title={Salida del Código #2},
  title after break={Salida del Código #2 (continuación)},
  listing file={#3},
  listing options={style=terminalstyle},
  colback=terminalBG,
  colupper=terminalFG,
  overlay
  first={
    % Close bottom (Red)
    \fill[btnClose] ([xshift=-15mm, yshift=-2.5mm]frame.north east) circle
    (3pt);
    % Minimize button (Yellow)
    \fill[btnMin] ([xshift=-10mm, yshift=-2.5mm]frame.north east) circle (3pt);
    % Maximize button
    \fill[btnMax] ([xshift=-5mm, yshift=-2.5mm]frame.north east) circle (3pt);
    \node[anchor=south east, font=\itshape\small, text=white!85] at (frame.south
    east) {(Continúa $\rightarrow$) \,};
  },
  overlay
  middle={
    \node[anchor=south east, font=\itshape\small, text=white!85] at (frame.south
    east) {(Continúa $\rightarrow$) \,};
  },
  % If the code box does not take up more than one page
  overlay
  unbroken={
    % Close bottom (Red)
    \fill[btnClose] ([xshift=-15mm, yshift=-2.5mm]frame.north east) circle
    (3pt);
    % Minimize button (Yellow)
    \fill[btnMin] ([xshift=-10mm, yshift=-2.5mm]frame.north east) circle (3pt);
    % Maximize button
    \fill[btnMax] ([xshift=-5mm, yshift=-2.5mm]frame.north east) circle (3pt);
  },
  #1
}

% ======== Example of use 
% \codeoutput[label=out:code2]{ref_for_title}{./path/to/output}

% ======== "Pseudocode" Language ======== %
\lstdefinelanguage{pseudocodigo}{
  morekeywords={
    Algoritmo,FinAlgoritmo, Si,FinSi,Mientras,FinMientras,
    Para,FinPara,Entonces,Hacer, Y,O,MOD,Escribir
  },
  sensitive=false
}

% ======== Pseudocode Style ======== %
\lstdefinestyle{pseudostyle}{
  language=pseudocodigo, basicstyle=\ttfamily\small,
  keywordstyle=\color{Psblue}\bfseries, commentstyle=\color{gray}\itshape,
  frame=none, numbers=left, breaklines=true, keepspaces=true
}
